\chapter{Abelian varieties of Weil type}\label{ch:weil}

In 1977 Weil wrote a short paper with a simple observation \cite{Weil77}.
Let an imaginary quadratic field act on an abelian variety so that its two
eigenspaces on the tangent space have the same dimension. Then the
determinant of the first cohomology over that field gives a plane of
Hodge classes, and for a general such variety these classes are not
products of divisor classes. This chapter proves Weil's observation,
describes the family in which these varieties live, computes the Hodge
ring of its general member, and shows that the Hodge conjecture for that
member is the algebraicity of one class. The last section looks at the
same classes on complex tori that are not algebraic, where Voisin showed
that they are not algebraic in any sense.

\section{Definition}

Fix an imaginary quadratic field $K=\QQ(\sqrt{-d})$, with $d$ a positive
squarefree integer, and write $\tau\mapsto\bar\tau$ for its nontrivial
automorphism. Fix one embedding of $K$ in $\CC$.

\begin{definition}\label{def:weiltype}
An abelian variety $A$ of dimension $2n$ with a polarisation $E$ is
\emph{of Weil type} for $K$ if it carries an embedding
$K\hookrightarrow\End^{0}(A)$ such that
\begin{enumerate}[label=(\roman*)]
\item the Rosati involution of $E$ restricts to complex conjugation on
$K$, and
\item the element $\sqrt{-d}$ acts on the tangent space at the origin with
the eigenvalues $\sqrt{-d}$ and $-\sqrt{-d}$ each of multiplicity $n$.
\end{enumerate}
\end{definition}

The second condition is the one that matters. Without it the
construction below still makes sense, but the classes it produces are not
of Hodge type.

Let $V=H^{1}(A,\QQ)$. Since $K$ is commutative, pullback makes $V$ a
$K$-vector space, of dimension $2n$ because $\dim_{\QQ}V=4n$. Over $\CC$
the action of $\sqrt{-d}$ splits $V_{\CC}$ into two eigenspaces
\[
  V_{\CC}=V_{+}\oplus V_{-},\qquad \dim V_{\pm}=2n,
\]
on which $\sqrt{-d}$ acts by $\pm\sqrt{-d}$. Complex conjugation of
coefficients exchanges them. Because $K$ acts by endomorphisms, it
preserves the Hodge decomposition, so each eigenspace splits as
$V_{\pm}=V_{\pm}^{1,0}\oplus V_{\pm}^{0,1}$. The space $V^{1,0}$ is dual to
the tangent space, so condition~(ii) says that $V^{1,0}_{+}$ and
$V^{1,0}_{-}$ both have dimension $n$. Since
$V^{0,1}_{+}=\overline{V^{1,0}_{-}}$, we get
\begin{equation}\label{eq:balanced}
  \dim V_{+}^{1,0}=\dim V_{+}^{0,1}=n,\qquad
  \dim V_{-}^{1,0}=\dim V_{-}^{0,1}=n .
\end{equation}

\subsection*{The hermitian form}
Condition~(i) says that $E(\tau x,y)=E(x,\bar\tau y)$ for $\tau\in K$. Fix
$\delta\in K$ with $\bar\delta=-\delta$, for instance $\delta=\sqrt{-d}$.
Then there is a unique $K$-hermitian form $H\colon V\times V\to K$, linear in
the first variable, with
\[
  E(x,y)=\Tr_{K/\QQ}\bigl(\delta\,H(x,y)\bigr).
\]
Indeed, for fixed $y$ the map $x\mapsto E(x,y)$ is $\QQ$-linear, and the
trace form identifies $\QQ$-linear forms on $V$ with $K$-linear forms; the
compatibility with $K$ makes the resulting form hermitian. Its
determinant, in a $K$-basis, is well defined in
$\QQ^{\times}/N_{K/\QQ}(K^{\times})$, and it is called the
\emph{discriminant}. The unitary group $\U(V,H)$ is the group of
$K$-linear automorphisms of $V$ that preserve $E$. Its subgroup of elements
of determinant one over $K$ is the special unitary group $\SU(V,H)$. Both
are algebraic groups over $\QQ$.

\begin{lemma}\label{lem:isotropic}
The eigenspaces $V_{+}$ and $V_{-}$ are isotropic for $E$, and $E$ pairs
them perfectly. So, as representations of $\U(V,H)_{\CC}\cong\GL(V_{+})$,
the space $V_{-}$ is the dual of $V_{+}$.
\end{lemma}

\begin{proof}
For $x,y\in V_{+}$ we have
$E(\sqrt{-d}x,\sqrt{-d}y)=E(x,\overline{\sqrt{-d}}\sqrt{-d}\,y)=dE(x,y)$, while
bilinearity gives $E(\sqrt{-d}x,\sqrt{-d}y)=-dE(x,y)$. So $E(x,y)=0$, and
the same holds on $V_{-}$. Since $E$ is nondegenerate on
$V_{\CC}=V_{+}\oplus V_{-}$ and vanishes on each summand, it pairs them
perfectly. An element of $\U(V,H)$ commutes with $K$, so it preserves
$V_{+}$ and $V_{-}$, and it is determined by its restriction to $V_{+}$; it
preserves $E$, so its action on $V_{-}$ is the contragredient.
\end{proof}

The form $h(x,y)=\delta E(x,\bar y)$, for $x,y\in V_{+}$, is hermitian on
$V_{+}$. It is the complexification of $H$ along the fixed embedding of
$K$, up to a positive constant. The Hodge--Riemann relations say that
$h$ is definite of one sign on $V_{+}^{1,0}$ and of the other on
$V_{+}^{0,1}$, so by \eqref{eq:balanced} it has signature $(n,n)$. This
is the form in which condition~(ii) is usually stated: the hermitian form
of an abelian variety of Weil type has signature $(n,n)$.

\section{The Weil classes}

\begin{definition}
The \emph{Weil plane} of $A$ is
\[
  \HW=\HW(A)=\Bigl(\textstyle\bigwedge^{2n}V_{+}\oplus\bigwedge^{2n}V_{-}\Bigr)
  \cap\textstyle\bigwedge^{2n}_{\QQ}V\ \subseteq\ H^{2n}(A,\QQ).
\]
Its elements are the \emph{Weil classes}.
\end{definition}

The two lines $\bigwedge^{2n}V_{\pm}$ sit inside
$\bigwedge^{2n}V_{\CC}=\bigoplus_{a+b=2n}\bigwedge^{a}V_{+}\otimes
\bigwedge^{b}V_{-}$ as the summands with $(a,b)=(2n,0)$ and $(0,2n)$.
Complex conjugation exchanges them, so their sum is defined over $\QQ$, and
$\HW$ is a rational plane. The field $K$ acts on $\HW$: an element $\tau$
acts on $\bigwedge^{2n}V_{+}$ by $\tau^{2n}$ and on $\bigwedge^{2n}V_{-}$ by
$\bar\tau^{2n}$. So $\HW$ is a one-dimensional $K$-vector space, canonically
the $K$-determinant $\bigwedge^{2n}_{K}V$.

\begin{theorem}[Weil]\label{thm:weilhodge}
The Weil classes are Hodge classes.
\end{theorem}

\begin{proof}
By \eqref{eq:balanced},
\[
  \textstyle\bigwedge^{2n}V_{+}=\bigwedge^{n}V_{+}^{1,0}\otimes
  \bigwedge^{n}V_{+}^{0,1},
\]
which is of type $(n,n)$. The same holds for $V_{-}$.
\end{proof}

The proof shows exactly why the balanced signature is needed. If the
eigenspace $V_{+}$ had $a$ holomorphic and $b$ antiholomorphic dimensions
with $a\ne b$, its top exterior power would have type $(a,b)$.

There is a second proof, through groups, which explains more. The Hodge
group of $A$ commutes with $K$ and preserves $E$, so it lies in
$\U(V,H)$. On $V_{+}$ an element $z$ of the unit circle acts through
$h(z)$ by $z^{-1}$ on $V_{+}^{1,0}$ and by $z$ on $V_{+}^{0,1}$, so its
determinant on $V_{+}$ is $z^{-n}z^{n}=1$. Thus $h(U(1))$ lies in the
special unitary group, and by minimality
\begin{equation}\label{eq:hginsu}
  \Hg(A)\subseteq\SU(V,H).
\end{equation}
The special unitary group acts trivially on $\bigwedge^{2n}V_{\pm}$,
since it acts there by the determinant. So the Weil classes are
invariants of $\Hg(A)$, which is another way of saying they are Hodge
classes.

The same computation shows why they are \emph{not} Lefschetz classes in
general. The centre of $\U(V,H)$ is the torus $\U(1)_{K}$ of norm-one
elements of $K$, and $\lambda\in\U(1)_{K}$ acts on $\bigwedge^{2n}V_{+}$ by
$\lambda^{2n}$.

\begin{proposition}\label{prop:weilnotlef}
Suppose $\End^{0}(A)=K$. Then the Lefschetz group of $A$ is $\U(V,H)$,
and no nonzero Weil class is a Lefschetz class.
\end{proposition}

\begin{proof}
The centraliser of $K$ in $\mathrm{Sp}(V,E)$ consists of the $K$-linear maps
that preserve $E$, which is $\U(V,H)$, a connected group. So
$L(A)=\U(V,H)$. Choose $\tau\in K^{\times}$ such that
$\lambda=\tau/\bar\tau$ is not a root of unity. This is possible because
$K$ contains finitely many roots of unity $\zeta$, and each equation
$\tau=\zeta\bar\tau$ cuts out a line in the plane $K$. Then $\lambda$ has
norm one. Then $\lambda^{2n}\ne1$ and $\bar\lambda^{2n}\ne1$,
so $\lambda$ moves every nonzero vector of $\HW$. By
Proposition~\ref{prop:lefinv}, Lefschetz classes are invariant under
$L(A)$, so they meet $\HW$ only in $0$.
\end{proof}

This is Weil's theorem: on an abelian variety of Weil type whose
endomorphism algebra is no bigger than $K$, the Weil classes are Hodge
classes that no polynomial in divisor classes can produce. When $n=1$ the
hypothesis never holds. The Weil classes then lie in $H^{2}$, so they are
divisor classes by the theorem of Lefschetz, and the proposition shows that
the endomorphism algebra of an abelian surface of Weil type is always
larger than $K$.

\section{The family}

All abelian varieties of Weil type for $K$ with a given hermitian space
$(V,H)$ fit into one family. We describe its base.

Fix $(V,H)$ of signature $(n,n)$, and let $h$ be the hermitian form on
$V_{+}$ as above. A polarised Hodge structure of weight one on $V$, of Weil
type for $K$, is determined by the $n$-dimensional subspace
$X=V_{+}^{1,0}\subset V_{+}$. Indeed, $V^{1,0}$ is isotropic for $E$ by
the Hodge--Riemann relations, so $V_{-}^{1,0}$ is the annihilator of $X$ in
$V_{-}$ under the pairing of Lemma~\ref{lem:isotropic}, and then
$V^{0,1}$ is the complex conjugate of $V^{1,0}$. The positivity part of
the Hodge--Riemann relations says that $h$, normalised with the right sign,
is positive definite on $X$. Conversely, every $n$-dimensional subspace on
which $h$ is positive definite arises in this way. So the parameter space
is
\[
  D=\{X\subset V_{+}:\dim X=n,\ h|_{X}>0\},
\]
an open subset of the Grassmannian of $n$-planes in a space of dimension
$2n$. It has dimension $n^{2}$, and it is the bounded symmetric domain of
the real group $\SU(n,n)$.

Let $\Gamma$ be a torsion-free arithmetic subgroup of $\SU(V,H)(\QQ)$
that preserves a lattice in $V$. The quotient
\[
  \mathrm{Sh}=\Gamma\backslash D
\]
is a smooth quasi-projective variety by the theorem of Baily and Borel
\cite{BB66}. Over it there is a polarised abelian scheme
$\pi\colon\mathcal A\to\mathrm{Sh}$ whose fibre over the class of $X$ is the
abelian variety with $H^{1}=V$ and $V_{+}^{1,0}=X$. We call it the
\emph{Weil family}. Its monodromy group on $H^{1}$ is $\Gamma$.

\begin{proposition}\label{prop:flat}
The Weil plane is a constant local subsystem of $R^{2n}\pi_{*}\QQ$, and
every Weil class extends to a global flat section that is a Hodge class at
every point.
\end{proposition}

\begin{proof}
The monodromy group $\Gamma$ lies in $\SU(V,H)$, which acts trivially on
$\HW$. So the monodromy fixes $\HW$ pointwise. By
Theorem~\ref{thm:weilhodge}, applied at each point, the transported
classes are Hodge classes everywhere.
\end{proof}

\begin{proposition}\label{prop:generalsu}
Let $n\ge2$. A very general member $A$ of the Weil family has
$\Hg(A)=\SU(V,H)$ and $\End^{0}(A)=K$. The set of members where the Hodge
group is smaller is a countable union of proper closed algebraic subsets
of $\mathrm{Sh}$.
\end{proposition}

\begin{proof}
The real group $\SU(n,n)$ has no compact factor, and $\SU(V,H)$ is simple
over $\QQ$, so by the Borel density theorem the arithmetic subgroup $\Gamma$
is Zariski dense in $\SU(V,H)$. By Andr\'e's theorem \cite{And92}, the
connected algebraic monodromy group is contained in the Hodge group of a
very general member. So that Hodge group contains $\SU(V,H)$, and by
\eqref{eq:hginsu} it equals it. For the endomorphisms, over $\CC$ the
group $\SU(V,H)$ is $\SL(V_{+})$, acting on $V_{+}$ by the standard
representation and on $V_{-}$ by its dual
(Lemma~\ref{lem:isotropic}). For $2n\ge3$ these are irreducible and not
isomorphic, so the commutant of $\SU(V,H)$ in $\End_{\QQ}(V)$ has
dimension two, and it contains $K$. By
Proposition~\ref{prop:hgabelian}(ii) it is $\End^{0}(A)$. For the last
statement, the members with a smaller Hodge group are those carrying an
extra Hodge class in some tensor space. Each such class has an algebraic
Hodge locus by Theorem~\ref{thm:cdk}, there are countably many rational
classes, and none of these loci is all of $\mathrm{Sh}$.
\end{proof}

We write $\mathrm{Sh}^{\circ}$ for the set of points where the Hodge group
is all of $\SU(V,H)$. Its complement is meagre.

\section{The Hodge ring of the general member}\label{sec:hodgering}

We can now compute every Hodge class on a member with full Hodge group.
The answer is as small as it could be, given Weil's theorem.

\begin{theorem}\label{thm:hodgering}
Let $n\ge2$ and let $A$ be a member of the Weil family with
$\Hg(A)=\SU(V,H)$. Then
\[
  \Hdg^{p}(A)=\QQ\,\eta^{p}\quad(p\ne n),\qquad
  \Hdg^{n}(A)=\QQ\,\eta^{n}\oplus\HW,
\]
where $\eta$ is the class of the polarisation.
\end{theorem}

\begin{proof}
By Proposition~\ref{prop:hgabelian}(iii) we need the invariants of
$\SU(V,H)$ in $\bigwedge^{*}V$. Over $\CC$,
\[
  \textstyle\bigwedge^{k}V_{\CC}=\bigoplus_{a+b=k}\bigwedge^{a}V_{+}\otimes
  \bigwedge^{b}V_{-},
\]
and as a representation of $\SL(V_{+})$ the summand is
$\bigwedge^{a}S\otimes(\bigwedge^{b}S)^{\vee}$, where $S=V_{+}$ is the
standard representation. Its invariants are the
$\SL(V_{+})$-homomorphisms $\bigwedge^{b}S\to\bigwedge^{a}S$. The exterior
powers $\bigwedge^{a}S$, $0\le a\le2n$, are irreducible, and two of them are
isomorphic only when $a=b$, or when both are trivial, which happens for
$a,b\in\{0,2n\}$. So there is a one-dimensional space of invariants in each
summand with $a=b$, and in the two summands $(a,b)=(2n,0)$ and $(0,2n)$,
and none elsewhere.

The summands with $(a,b)=(2n,0)$ and $(0,2n)$ make up the Weil plane. In
the summand with $a=b=p$, the class $\eta^{p}$ is a nonzero invariant for
$0\le p\le2n$. Indeed $\eta$ is invariant, and $\eta^{p}\ne0$ because
$\eta^{2n}$ is a positive multiple of the volume class. Its component in
that summand is nonzero, because the other summands of $\bigwedge^{2p}V$
contain no invariants except the Weil plane when $p=n$, and $\eta^{n}$ is
not a Weil class: the centre $\U(1)_{K}$ fixes $\eta^{n}$ and moves every
nonzero Weil class. So $\eta^{p}$ spans the invariants with $a=b=p$.
\end{proof}

\begin{corollary}\label{cor:oneclass}
Let $n\ge2$ and $A\in\mathrm{Sh}^{\circ}$. The following are equivalent.
\begin{enumerate}[label=(\roman*)]
\item The Hodge conjecture holds for $A$ in every degree.
\item Some nonzero Weil class on $A$ is algebraic.
\item Every Weil class on $A$ is algebraic.
\end{enumerate}
\end{corollary}

\begin{proof}
By Theorem~\ref{thm:hodgering}, (i) and (iii) are equivalent, since the
powers of $\eta$ are algebraic. Suppose a nonzero Weil class $\omega$ is the
class of a cycle $Z$. Choose $\tau\in K$ with $\tau^{2n}\notin\QQ$, for
example $\tau=a+\sqrt{-d}$ with $a$ a positive integer chosen so that
$\tau/\bar\tau$ is not a root of unity. Some positive integer multiple $m\tau$
is an endomorphism of $A$, and its pullback acts on $\HW$ by
$(m\tau)^{2n}$. So $(m\tau)^{*}Z$ is a cycle with class $m^{2n}\tau^{2n}\omega$.
The classes $\omega$ and $\tau^{2n}\omega$ are linearly independent over
$\QQ$, since $\HW$ is a one-dimensional $K$-space and $\tau^{2n}\notin\QQ$.
So they span $\HW$, and every Weil class is algebraic.
\end{proof}

For a very general abelian variety of Weil type, then, the whole Hodge
conjecture is a statement about one class. That is the sense in which these
varieties are the first test case. They are also, by
Proposition~\ref{prop:weilnotlef}, exactly the case that divisors cannot
reach.

\section{Weil tori}\label{sec:voisintori}

Drop the polarisation, and keep only the action of $K$ with balanced
signature. Let $V$ be as before. For complementary $n$-dimensional
subspaces $X,Y\subset V_{+}$, the subspace $X\oplus\overline{Y}$ of
$V_{\CC}$ is the space of holomorphic one-forms of a complex structure on
the real torus $V_{\RR}^{\vee}/H_{1}$. Call the resulting complex torus
$T_{X,Y}$. On it $K$ acts with balanced signature, and every complex
structure commuting with $K$ with balanced signature arises in this way.
The pairs $(X,Y)$ form a connected complex manifold $\mathcal S$ of
dimension $2n^{2}$, and the Weil family is the submanifold of dimension
$n^{2}$ where $Y$ is the $h$-orthogonal complement of $X$ and $h$ is
positive on $X$. Since
$\bigwedge^{2n}V_{+}=\bigwedge^{n}X\otimes\bigwedge^{n}Y$, with $X$
holomorphic and $Y$ antiholomorphic, the Weil classes are of type $(n,n)$
on every $T_{X,Y}$. These are the \emph{Weil tori}.

\begin{proposition}\label{prop:weiltorushodge}
For a very general Weil torus $T$, the Hodge classes are $\HW$ in degree
$2n$ and zero in every other degree strictly between $0$ and $4n$.
\end{proposition}

\begin{proof}
There are countably many rational classes, and the Hodge locus of each is
an analytic subset of $\mathcal S$. So a class that is a Hodge class at a very
general point is a Hodge class at every point. Let
$G=\operatorname{Res}_{K/\QQ}\GL_{K}(V)$ be the group of $K$-linear
automorphisms of $V$. Its real points, which form the group $\GL(V_{+})$,
act transitively on pairs of complementary subspaces of $V_{+}$, hence on
$\mathcal S$. The Mumford--Tate group $M$ of a very general member lies in $G$,
because $K$ acts by endomorphisms of every member, and it contains the
Mumford--Tate group of every member. So the real Lie algebra of $M$
contains the complex structure of every member, and hence the span of the
orbit of one of them under the adjoint action of $G(\RR)$. That span is an
ideal of $\mathfrak{gl}(V_{+})$. It contains a non-central element, the
complex structure, which acts by $i$ on $X$ and by $-i$ on $Y$. So it
contains $\mathfrak{sl}(V_{+})$, and $M$ contains the derived group of $G$.
Over $\CC$ that derived group is $\SL(V_{+})\times\SL(V_{-})$, acting
independently on the two factors of
$\bigwedge^{*}V_{+}\otimes\bigwedge^{*}V_{-}$. Its invariants are spanned by
$1$, $\bigwedge^{2n}V_{+}$, $\bigwedge^{2n}V_{-}$ and their product. Hodge
classes are invariant under $M$, which proves the statement.
\end{proof}

Notice the difference from the polarised case. There the two factors are
dual to each other under one group, and the classes $\eta^{p}$ appear. On a
very general Weil torus there is no polarisation, and nothing but the Weil
plane survives.

\begin{lemma}\label{lem:weiltorussub}
A very general Weil torus has no analytic subsets other than finite sets
and itself.
\end{lemma}

\begin{proof}
Let $Z$ be an irreducible analytic subset of codimension $c$ with
$0<c<2n$. For a K\"ahler class $\kappa$ we have
$\int_{T}[Z]\cup\kappa^{2n-c}=\int_{Z}\kappa^{2n-c}>0$, so $[Z]$ is a nonzero
Hodge class of degree $2c$. By Proposition~\ref{prop:weiltorushodge}, $c=n$
and $[Z]\in\HW$. Take for $\kappa_{0}$ the class of a constant positive
definite hermitian form that is block diagonal for the splitting
$H^{1,0}=X\oplus\overline{Y}$. Its two blocks lie in $X\otimes\overline{X}$
and $\overline{Y}\otimes Y$, both inside $V_{+}\otimes V_{-}$. So every term
of $\kappa_{0}^{n}$ has $n$ factors from $V_{+}$, and its product with an
element of $\bigwedge^{2n}V_{+}$ vanishes, since $\bigwedge^{3n}V_{+}=0$. The
same holds with $V_{-}$. Hence $\int_{T}[Z]\cup\kappa_{0}^{n}=0$, a
contradiction.
\end{proof}

Voisin proved much more. Coherent sheaves are the natural generalisation of
subvarieties on a complex manifold that need not be projective, and one can
ask whether the Weil classes are at least combinations of Chern classes of
coherent sheaves. They are not.

\begin{theorem}[Voisin \cite{Voi02b}]\label{thm:voisin}
Let $n=2$ and let $T$ be a very general Weil torus. Every coherent analytic
sheaf on $T$ has $c_{2}=0$ in $H^{4}(T,\QQ)$. In particular the Weil classes
of $T$ are Hodge classes that are not combinations of Chern classes of
coherent sheaves.
\end{theorem}

We sketch the proof, because the same argument works for every $n\ge2$ and
it is used in Chapter~\ref{ch:remaining}. Chern classes of coherent sheaves
on a compact complex manifold exist in rational Deligne cohomology
\cite{Gri10}, and they are Hodge classes. By Lemma~\ref{lem:weiltorussub}, the torsion of a
coherent sheaf and the locus where a torsion-free sheaf fails to be
reflexive are supported on finite sets, and they contribute only in top
degree. So it is enough to treat reflexive sheaves. Since there are no
Hodge classes in degree two, every sheaf has $c_{1}=0$, so every
torsion-free sheaf has slope zero for every K\"ahler class. Stability then
does not depend on the K\"ahler class, and a Jordan--H\"older filtration
reduces the question to stable reflexive sheaves $Q$. For these, the
Bogomolov inequality, in the form of Bando and Siu for reflexive sheaves
\cite{BS94}, gives
\[
  \int_{T}c_{2}(Q)\cup\kappa^{2n-2}\ge0
\]
for every K\"ahler class $\kappa$, with equality only if $Q$ is a flat
vector bundle. For $n\ge3$ the class $c_{2}(Q)$ is a Hodge class of degree
four, hence zero. For $n=2$ it lies in $\HW$, and the class $\kappa_{0}$ of
the proof of Lemma~\ref{lem:weiltorussub} gives equality. Either way $Q$ is
flat, and its Chern classes vanish. The same argument shows that every
coherent sheaf on a very general Weil torus of dimension $2n$ has
$\operatorname{ch}_{k}=0$ for $0<k<2n$.

Voisin's theorem is a counterexample to the Hodge conjecture for compact
K\"ahler manifolds, as explained in Chapter~\ref{ch:conjecture}. For us it
plays a different role. It is a warning about deformation arguments. The
Weil class is of Hodge type on every Weil torus, polarised or not. An
argument that moves a cycle along every deformation preserving the Hodge
type of the Weil class would move it onto a Weil torus, where no cycle or
sheaf represents the class. So any such argument has to use the
polarisation in an essential way. This is the source of the main obstruction
in Chapter~\ref{ch:remaining}.
